\documentclass[11pt,oneside]{report}
\usepackage[T1]{fontenc}
\usepackage{lmodern}
\usepackage[utf8]{inputenc}
\usepackage[english]{babel}

\usepackage{amsmath,amssymb,amsthm}
\usepackage{booktabs}
\usepackage{longtable}
\usepackage{array}
\usepackage{listings}
\usepackage{xcolor}
\usepackage[hidelinks]{hyperref}
\usepackage[capitalise,noabbrev]{cleveref}
\usepackage{microtype}
\usepackage{enumitem}
\usepackage{geometry}
\geometry{a4paper,margin=1.1in}

% --- notation ---------------------------------------------------------------
\newcommand{\gid}{\ensuremath{g}}
\newcommand{\iv}{\ensuremath{\mathbf{i}}}
\newcommand{\lev}{\ensuremath{\ell}}
\newcommand{\dvol}{\ensuremath{\Delta V}}
\newcommand{\dsp}{\ensuremath{d}}                 % spatial dimension
\newcommand{\sublev}[1]{\ensuremath{\{f > #1\}}}
\newcommand{\betti}[1]{\ensuremath{b_{#1}}}
\newcommand{\eulerchi}{\ensuremath{\chi}}
\newcommand{\Tor}[1]{\ensuremath{\mathbb{T}^{#1}}}

% code listings: C++ and shell, both in the same restrained style
\definecolor{codebg}{gray}{0.96}
\definecolor{codecmt}{rgb}{0.30,0.42,0.30}
\definecolor{codekw}{rgb}{0.20,0.20,0.55}
\lstdefinestyle{plain}{
  basicstyle=\ttfamily\footnotesize,
  keywordstyle=\color{codekw},
  commentstyle=\itshape\color{codecmt},
  stringstyle=\color{black},
  backgroundcolor=\color{codebg},
  frame=single, framerule=0pt, framesep=4pt,
  breaklines=true, breakatwhitespace=true,
  showstringspaces=false, columns=fullflexible,
  captionpos=b, aboveskip=1.0em, belowskip=1.0em,
}
\lstdefinestyle{cpp}{style=plain,language=C++}
\lstdefinestyle{sh}{style=plain,language=bash}
\lstset{style=plain}

% a lightweight environment for the recurring "rule learned from a defect"
\newtheoremstyle{ruleplain}{1em}{1em}{\itshape}{}{\bfseries}{.}{ }{}
\theoremstyle{ruleplain}
\newtheorem{designrule}{Rule}[section]

\newcommand{\decide}[1]{\medskip\noindent\textbf{Decision required.} #1\medskip}


\title{\bfseries A Conforming Multilevel Complex for Topological Analysis of
       AMReX Plotfile Data\\[0.4em]
       \large Design, Verification, and Application to Cellular Structure in
       Lean Premixed Hydrogen Flames}
\author{Marc Day}
\date{\today}

\begin{document}
\maketitle

\begin{abstract}
\noindent
This note records the design and verification of \texttt{AMRComplex} and its
consumers in \texttt{PeleAnalysis}: a cell complex built over an AMReX plotfile
hierarchy that is conforming across coarse--fine interfaces and correct on
periodic domains, together with marching-simplex contouring, threshold
segmentation, and the test apparatus that establishes both.

The motivating problem is the characterisation of cellular structure produced
by the thermodiffusive instability in lean premixed hydrogen flames, where the
definition of a ``cell'' is itself unsettled. That shapes the design: the
tooling exists to \emph{compare} candidate definitions within one framework
rather than to compute a single privileged one.

Two themes recur and are recorded deliberately. First, a single design
choice---that a node's identity is its \emph{index}, not its
\emph{position}---removes an entire class of defect, and most of what follows
is a consequence of it. Second, several cheap and plausible verification
checks were found to pass on demonstrably broken output; the negative controls
that expose them are retained in the source and are documented here as part of
the design rather than as an afterthought.
\end{abstract}

\tableofcontents

\chapter{The problem, and what it demands of the tooling}
\label{ch:problem}

\section{Physical target}

The driver is the characterisation of the effects and time-dependent dynamics
of the thermodiffusive instability in lean premixed hydrogen flames. These
flames spontaneously form cellular structures whose sizes and separations
change dramatically with fuel mixture and with the turbulence level of the
approach flow. The analysis lineage is that of Bremer
et~al.~\cite{bremer2010burning}; the present work is a rebuild of the
underlying geometric and topological machinery on modern AMReX.

Relevant configurations include a doubly periodic \emph{flamesheet} with inflow
and outflow in the third direction, and growing flame kernels, each in two and
three dimensions, and computed with different models for the transport physics
that drives the instability.

\section{What is not known}

The tooling exists because a prior question is open: \emph{what a cell is}.

There are chemistry arguments about when a local flame is self-sustaining, and
there is the distinct notion of \emph{back-supported} structure. Separately,
observation supplies a set of surrogates that have traditionally stood in for
the presence of a flame---fuel consumption rate, temperature above some
threshold, a flame intermediate, and others. Which surrogate to use, and at
what threshold, is not settled, and the answer may not be the same for all of
them.

A complication specific to these flames sharpens the difficulty. Because they
are far from unity Lewis number, the progress variables do not align: contours
of fuel, of temperature and of intermediates cross one another. Each may
therefore tell a different story in detail while supporting a similar global
picture---and whether that is so is itself one of the questions to be answered.

\section{Requirements that follow}

\begin{enumerate}[label=(R\arabic*),leftmargin=*]
\item\label{req:compare} \textbf{Compare, do not privilege.} The deliverable is
  not the topology of a chosen field but a framework in which candidate
  definitions are computed identically and compared directly. Adding a
  candidate surrogate should be a command-line argument, not a code change.
\item\label{req:robust} \textbf{Robustness to definition.} Any reported measure
  must come with a statement of how sensitive it is to the threshold and field
  that produced it. A number without that statement is not usable here.
\item\label{req:periodic} \textbf{Periodic correctness with arbitrary face
  placement.} Doubly and triply periodic domains are routine, and the location
  of the domain faces carries no physical meaning. No reported quantity may
  depend on where they happen to fall.
\item\label{req:amr} \textbf{Correctness across coarse--fine interfaces.} The
  data are adaptively refined; an analysis that is only approximately right at
  level transitions is not acceptable, because the structures of interest are
  exactly where the refinement is.
\item\label{req:parallel} \textbf{Parallel throughout.} Snapshots are large.
  Any step that gathers the whole problem to one rank is a design fault, not an
  optimisation opportunity.
\item\label{req:track} \textbf{Relate structures across snapshots.} The
  dynamics are the point, so features must be identifiable from one time to the
  next.
\item\label{req:sample} \textbf{Carry sampled data onto extracted geometry.}
  Surfaces must be able to carry an arbitrary vector of interpolated field
  values at their nodes; this is what makes conditional analysis possible.
\end{enumerate}

\Cref{req:robust} deserves emphasis because it determines the topological
target in \cref{ch:topology}: it is a request for \emph{persistence}, and
persistence is what a merge tree supplies.

\chapter{Assessment of the existing tools}
\label{ch:assessment}

Two existing tools in \texttt{PeleAnalysis} bear on this work.
\texttt{isosurface.cpp} extracts a multilevel isosurface by marching cubes;
\texttt{amrToFE.cpp} converts an AMR hierarchy into a finite-element mesh. The
first is the natural starting point and turns out to be the wrong one; the
second contains the construction that is actually needed.

\section{\texttt{isosurface.cpp}: a good front end, the wrong engine}

The plotfile front end is directly reusable and is the tedious part to get
right: \texttt{PlotFileData} access, component selection by name, per-level
\texttt{Geometry} carrying an explicit periodicity, \texttt{FillPatchTwoLevels}
for level-to-level data, and the fine-covered mask with its periodic-shift
pass. Roughly the first 250 lines of \texttt{main()} survive.

The engine does not, for reasons that are structural rather than qualitative.

\subsection{One isovalue is not a field topology}

Marching cubes yields a triangle soup at a single isovalue with no domain
connectivity attached. Every contour-tree and Reeb-graph
algorithm~\cite{carr2003contourtree,tierny2009loopsurgery,gueunet2016contourforests}
consumes a cell complex with vertex scalars and explicit adjacency, sweeping
sorted vertices with union--find. \texttt{Polygonise} discards precisely that:
its output is triangles of iterators into a map keyed on edges. There is no
path from that representation to critical points and their pairing.

\subsection{The coarse--fine seam is only approximately closed}

Coarse dual cells touching a covered cell emit nothing, and the fine level
covers them through its ghost band. Coarse and fine nodes on the interface
plane then land at different positions and are welded only if they compare
close. For an area integral the resulting cracks are a percent-level error; for
topology they are fatal, since a crack splits one component into two and
invents both a feature and a critical point.

\subsection{Node welding is not a strict weak ordering}

The node comparator short-circuits on an $\epsilon$ test before falling through
to a lexicographic comparison. With $a \sim b$ and $b \sim c$ but
$a \not\sim c$, transitivity fails and the behaviour of the ordered container
becomes insertion-order dependent. At $\epsilon = 10^{-15}$ this is effectively
exact comparison and largely harmless; any attempt to widen the tolerance in
order to weld coarse nodes to fine ones converts it into non-deterministic
topology.

\subsection{Piecewise-constant interpolation in the ghost band}

The coarse-to-fine fill is piecewise constant, so the field a fine box sees in
its ghost region is a staircase of coarse values. The two sides of the
interface therefore see genuinely different functions, and the staircase
manufactures exactly-flat plateaus, which are degenerate for any sweep and
create spurious critical points at every interface.

\subsection{Collation to a single rank}

Nodes and elements are gathered to the I/O processor, which rebuilds the global
sets serially. This already binds for isosurfaces at scale and is disqualifying
under~\ref{req:parallel}.

\subsection{The periodic defect is structural}

The tool's own diagnostics record it: across a periodic face the crossing is
interpolated from raw coordinates rather than the minimum image, so the node
lands mid-domain---observed inflating a reported arc length by 45\%\ in one
case. The cause is that connectivity is inferred \emph{from coordinates}. In a
complex whose adjacency is topological, the defect cannot be expressed.

\section{\texttt{amrToFE.cpp}: the construction that was needed}

This tool builds a conforming multilevel mesh, and the reason it works is a
single choice: a node's identity is the pair $(\lev,\iv)$ of level and index,
never a coordinate. Two elements share a vertex because they \emph{name the
same cell}.

The coarse--fine transition follows from that. For $\lev > 0$ the ring of cells
within \texttt{ref} of the fine grids but outside them is stamped with the
\emph{coarse parent's node}: for each coarse cell $\iv$, every one of the
$r^{\dsp}$ fine positions $r\iv + \iv_t$ receives the node
$(\lev-1, \iv)$. Coarse cells under a finer grid are marked covered, and an
element is emitted only when all of its corners are valid. Interface elements
are therefore emitted at the fine level with their coarse-side corners being
the coarse level's own nodes, by identity. Because $r^{\dsp}$ fine positions
collapse onto one coarse node, those elements carry repeated indices: they are
degenerate hexahedra---pyramids, prisms and tetrahedra written in brick
form---and that degeneracy \emph{is} the stitching. It is an alternative to the
explicit stitch-cell case tables of Weber
et~al.~\cite{weber2001crackfree,weber2012parallel} that needs no case table and
imposes no restriction on refinement ratio.

What it lacked for the present purpose:

\begin{itemize}[leftmargin=*]
\item \textbf{Periodicity by duplication, not identification.} Its periodic
  option grows the domain past the face and appends shifted copies of the
  boxes. The stamped indices there lie \emph{outside} the domain and are
  therefore different nodes from their partners on the far side. The result is
  an extended slab, correct for a seamless picture and wrong for topology: the
  complex is not a torus, so it carries no $1$-cycles.
\item \textbf{No parallelism.} The node map is built per rank with a per-rank
  counter, so identifiers collide between ranks, and the gather is commented
  out. The tool is single-rank-correct only, and the serialness is in the data
  structure rather than the interface.
\item \textbf{Legacy interfaces} predating the modern AMReX API.
\end{itemize}

\Cref{ch:complex} describes the rebuild that keeps the identity scheme and
addresses each of these.

\chapter{The complex}
\label{ch:complex}

\texttt{AMRComplex} builds a single cell complex over the whole plotfile
hierarchy. Everything in this chapter follows from one decision, stated first.

\section{Identity is an index, not a position}
\label{sec:identity}

A node is the pair $(\lev,\iv)$: the level and the integer index of the cell
whose centre it occupies. It is stored as a plain aggregate so that it is
trivially copyable and can therefore live inside a \texttt{BaseFab} and travel
through \texttt{FillBoundary} and \texttt{ParallelCopy}.

\begin{designrule}\label{rule:identity}
Two elements share a vertex because they name the same cell, never because two
coordinates compared close.
\end{designrule}

The ordering on nodes is lexicographic---by level, then by index---which is a
genuine strict weak ordering, unlike the $\epsilon$ comparison of
\cref{ch:assessment}. Three substantial consequences follow, and they are the
substance of the design rather than incidental benefits.

\section{Periodic identification is free}
\label{sec:periodic}

\texttt{FillBoundary} copies node \emph{values} into ghost cells. It cannot
shift an index stored inside the object it copies. Therefore, after a periodic
fill, the ghost beyond the high face literally contains the low-face node. That
copy \emph{is} the periodic identification: there is no arithmetic, no modular
reduction, and no opportunity for the coordinate-based defect of
\cref{ch:assessment}.

This satisfies~\ref{req:periodic} exactly. No boundary plane enters any
calculation, so the arbitrary placement of the domain faces cannot affect any
reported quantity---which is why the doubly and triply periodic cases cost
nothing here, a point that returns decisively in \cref{ch:topology}.

\section{Coarse--fine stitching}
\label{sec:stitch}

The construction is that of \cref{ch:assessment}: fine ghosts lying outside the
level are stamped with their coarse parent's node, so that a block of
$r^{\dsp}$ fine positions resolves to one coarse node and interface elements
become degenerate cells. Refinement ratio is unrestricted; ratios $2$ and $4$
are both verified.

\paragraph{Proper nesting is a precondition, not a nicety.}
A level's ghost ring must find valid parents one level down. If the parent level
is itself covered there by this level, the stitch has nothing to attach to and
the interface silently leaks volume. The first three-level test failed with a
$3\%$ residual for exactly this reason---a defect in the test generator, but one
that nothing detected. Construction now aborts when a ghost stamped from level
$\lev-1$ comes back covered, naming the levels involved.

\section{Element ownership}
\label{sec:ownership}

Elements must be generated exactly once, globally, without a deduplication
pass. The obvious rule is wrong, and wrong in a way that silently removes
exactly the elements the design exists to produce.

Consider a one-dimensional slice with refinement ratio $2$, a fine region
occupying coarse cells $[8,24)$ and therefore fine indices $[16,48)$. Coarse
cell $7$ is valid, at $7.5\,\dvol_c^{1/\dsp}$; coarse cell $8$ is covered; the
first fine cell centre is at $8.25$ in coarse units. The bridging element spans
$[7.5, 8.25]$ and its base cell is fine index $15$---a \emph{ghost}. No fine
box holds it as valid, and the coarse box that would hold it has a covered
corner and is rejected. Under the rule ``the box whose valid region contains the
base cell'', nobody generates it, leaving a gap one stitch-layer thick around
every refined region. Measured on a two-level $32^3$ case: \textbf{1.93\%\ of
the domain volume}.

The rule that works: the base range is the valid box grown by one \emph{on the
low side}, and a base cell in that shell belongs to the lowest-indexed box that
would visit it. Since a \texttt{BoxArray} is globally replicated this is
deterministic and requires no communication: the same element is produced by
exactly one rank at any rank count. Periodic wrap images are excluded, since the
identical element---the same node set, identity being topological---is already
produced where the cell is valid.

Ownership is resolved as box algebra once per grid, not as a query per cell: the
low-side shell is at most $\dsp$ boxes, the valid-elsewhere and periodic-image
exclusions are \texttt{BoxList} subtractions, and the lower-index claim is one
\texttt{intersections} call.

\section{Global numbering}
\label{sec:numbering}

A \texttt{BoxArray} is a globally replicated object: \texttt{ba[i]} is the same
box on every rank for any number of ranks, and only the
\texttt{DistributionMapping} depends on the rank count. Its own index order is
therefore already a deterministic global ordering, and
\begin{equation}
  \gid(\lev,\iv) \;=\; O_{\lev} \;+\; o_{\lev,i} \;+\; \mathrm{lex}(\iv, B_{\lev,i}),
  \label{eq:gid}
\end{equation}
where $O_{\lev}$ is the cumulative cell count of levels below $\lev$,
$o_{\lev,i}$ the cumulative count of boxes below $i$ within level $\lev$, and
$\mathrm{lex}$ the lexicographic offset within box $B_{\lev,i}$.

\Cref{eq:gid} needs \emph{no communication at all}: every term is a box size,
known globally. It counts every cell rather than only the uncovered ones, which
costs nothing---no array is allocated per identifier---and buys a great deal:
$\gid$ becomes a pure function of $(\lev,\iv)$, so it can be assigned at fill
time and then travels with the identity through \texttt{FillBoundary}, with no
second pass to propagate identifiers into ghosts. Identifiers are consequently
sparse over valid nodes; a compaction pass can be added for formats that need
them contiguous.

\section{Field data}
\label{sec:fields}

An arbitrary vector of fields is carried on the nodes, on the same grids with
the same ghost width as the identity, and propagated by the same calls: one
\texttt{FillBoundary} with the same periodicity, and the coarse-parent stamp
performing identity and value \emph{in the same loop under the same test}.

Stamping them in a second loop would be a second opportunity for a stitch
node's value to disagree with the identity it belongs to, and nothing
downstream would notice: the geometry would still verify and the extracted
surface would simply be wrong. Fields are initialised to a sentinel so that an
unpropagated value is detectable rather than plausible, and the verification
counts valid nodes still carrying it.

\section{The canonical traversals}

Two template traversals are the only sanctioned access:
\texttt{forEachElement}, which yields each element's corner nodes and their
field values, and \texttt{forEachNode}, which yields each valid node once with
its cell volume and a flag marking nodes another rank may also see. Every
consumer---verification, contouring, segmentation---goes through them, so that
ownership, corner ordering and the validity test are stated once. Two copies of
such a loop are two opportunities to disagree about which elements exist.

\chapter{Contouring}
\label{ch:contour}

\texttt{AMRContour} extracts an isosurface in three dimensions and a contour
curve in two, from the complex rather than from the grid.

\section{Marching simplices, not marching cubes}

The complex is subdivided into simplices---tetrahedra in 3D, triangles in
2D---and contoured with marching tetrahedra (respectively triangles). There are
three distinct cases and no ambiguous faces, so the extracted surface is
topologically consistent by construction. There is no asymptotic-decider
question to answer~\cite{chernyaev1995mc33} and none of the several hundred
lines of case tables that \texttt{isosurface.cpp} carries.

\section{Subdivision: cone from the minimum identifier}
\label{sec:cone}

Each cell is coned from its \emph{minimum-$\gid$} corner. In 3D this means the
three faces not containing the apex, each split on its own minimum-$\gid$
diagonal; in 2D, the two edges not containing it.

Every face of the cell then ends up split by the diagonal through its own
minimum-$\gid$ corner. For the faces containing the apex this is automatic,
since a corner minimal over the cell is minimal over any face containing it;
for the opposite faces it is imposed. Two cells sharing a face see the same
corner identifiers and therefore choose the same diagonal.

\begin{designrule}\label{rule:noparity}
Do not key the diagonal choice on index parity. Parity differs between levels,
so a parity rule fails exactly where two levels meet---the one place the
subdivision must be right.
\end{designrule}

This construction was adopted only after establishing the premise it depends
on. That two cells at a coarse--fine interface share a face \emph{as a matching
set of nodes} is not obvious: the stitch cells' coarse-side faces collapse to
points or edges for most index parities, and a T-junctioned complex would tile
the domain by volume while failing to share faces. The face-incidence census of
\cref{sec:faces} settles it by measurement.

Degenerate simplices---those with repeated identifiers, arising from stitch
cells---are dropped, which is what turns a degenerate hexahedron into the
pyramid or prism it actually is.

\section{Surface vertices are named by their edge}
\label{sec:edgekey}

A surface vertex is identified by the edge it lies on, and that edge is named by
the unordered pair of node identifiers at its ends. Two elements sharing an
edge therefore find the same vertex by identity, at any level and across the
periodic wrap---no coordinate welding, no tolerance, and none of the crack and
mis-placement failures that follow from either.

Everything about a vertex must be a function of the edge alone. This was not
initially true: the position was reconstructed as \emph{origin plus offset} in
the frame of the enclosing element, and two elements sharing an edge have
different origins, so they produced copies differing by an ulp. Within a rank
the cache hid this completely---only the first element to reach an edge ever
computed it---and it surfaced only as two ranks disagreeing. The position is now
computed from the two end nodes' own coordinates, with the far node brought to
its minimum image of the near one:
\begin{equation}
  \mathbf{x}_v \;=\; \mathbf{x}_0 + t\,\bigl(\mathbf{x}_1^{\,\mathrm{mi}} - \mathbf{x}_0\bigr),
  \qquad
  t = \frac{v - f_0}{f_1 - f_0},
  \label{eq:vertpos}
\end{equation}
with the pair ordered by identifier so that $t$ is likewise independent of who
asked.

\begin{designrule}\label{rule:edgefn}
Nothing about the element that reached an edge may enter what its vertex is.
\end{designrule}

\section{Orientation}
\label{sec:orient}

Nothing upstream constrains the winding: vertex order emerges from the marching
case, which follows the simplex's corner order, which follows the cone
construction---all arbitrary relative to the geometry. Measured on a three-level
sphere before this was imposed: \textbf{64598 of 157680 boundary items
misoriented}, about $41\%$, which is what random produces. The corresponding
check in \texttt{isosurface.cpp} is an empty function body.

Each element is now oriented so its normal points toward \emph{increasing} $f$.
The reference direction is from the below-value corners of the simplex to its
above-value corners: exact, local, and independent of the simplex's handedness.

Verification is combinatorial rather than a comparison of normals. A closed
surface is consistently oriented exactly when every interior edge is traversed
in \emph{opposite} directions by the two elements sharing it. That is reported
two ways: a rank-local exact count, and a signed modular sum in which each
traversal contributes $\pm$ its hash, so a correct pair cancels wherever the two
elements live and the global result is exactly zero at any rank count. In 2D the
same statement applies to segment endpoints---each vertex the head of one
segment and the tail of the other.

As an independent geometric confirmation, the enclosed volume by the divergence
theorem on a sphere of radius $0.3$ evaluates to $+0.1130050$ against an exact
$0.1130974$: positive to $0.08\%$, so the normals point outward rather than
merely agreeing with one another.

\section{Carried components}
\label{sec:carried}

Requirement~\ref{req:sample}. Every carried field is interpolated at the same
parameter $t$ of \cref{eq:vertpos} that the isosurface component determines,
which is the convention of \texttt{isosurface.cpp}. Because the vertex is cached
on its edge key, a carried value is computed \emph{once} and every element
sharing that edge receives the identical result---across levels and across the
wrap---rather than each recomputing and the copies being reconciled afterwards
by comparing coordinates.

One caveat is worth recording. At a stitch edge the two endpoints lie on
different levels, so a carried value there is linearly blended across a
resolution jump over a physically longer edge. That is the truthful thing to do
with the data available, and it applies equally to the surface position itself,
but sampled values in the stitch band are a blend across levels rather than a
fine-level quantity.

\section{Output}

MEF is written in the format the existing \texttt{*MEF} chain reads: title,
variable names, element and node-per-element counts, a node block as an
\texttt{FArrayBox} in node-major order, and one-based connectivity. Under MPI it
gathers to the I/O processor---the same wall \texttt{isosurface.cpp} has, and the
reason a \texttt{.pvtu} path remains on the roadmap. It differs in one respect
that matters: duplicate vertices merge by \emph{edge key} rather than by
coordinate proximity, so the merge is exact and the global node ordering, a sort
on those keys, is identical at any rank count.

Both orderings must be canonicalised. Sorting the nodes and not the elements
left the node block byte-identical across rank counts and the connectivity block
differing from its first byte, which is how the defect was localised without
further instrumentation.

\chapter{Segmentation and the threshold sweep}
\label{ch:segment}

\texttt{AMRSegment} finds the connected components of a threshold set, and
\texttt{amrSweep} walks the threshold and reports how they behave. Together they
are the substrate for \ref{req:compare}: the smallest apparatus that produces a
scientific result about candidate cell definitions.

\section{Components of a threshold set}

Components of $\sublev{v}$---or of $\{f < v\}$---are found by union--find over
the complex's own adjacency, which already carries the periodic wrap and the
coarse--fine stitch. No boundary plane enters, so \ref{req:periodic} holds with
nothing further required.

\paragraph{Labels are the minimum identifier in the component.}
This is a property of the component itself and not of who found it, so the same
run at any rank count produces identical labels. Segmentations are therefore
directly comparable across rank counts, across time snapshots, and across
different candidate fields---which is precisely what \ref{req:compare}
and~\ref{req:track} require.

\paragraph{The cross-rank merge exploits the same fact.}
Each rank publishes $(\gid, \text{local root})$ pairs; two ranks naming the same
$\gid$ with different roots is exactly the information needed to join their
components, and every rank replays the same union--find over the gathered pairs,
so no broadcast is needed.

Only nodes another rank might also see are published: ghosts, owned nodes within
one cell of a box boundary, and---the case easy to miss---owned nodes adjacent to
a covered cell, since those may be stamped as coarse parents into a finer box's
ghost region and so be seen by whichever rank owns that fine box. Omitting the
third case would fail to join components across a coarse--fine interface,
silently and only under MPI. Publishing every node instead, as the first version
did, is $O(N)$ all-to-all: it passes at four ranks on a $64^3$ box and fails on
the snapshots the tool exists for.

\section{Volumes are integer-derived}
\label{sec:intvol}

Cell volumes take exactly one value per level. The tally therefore counts
integers per (component, level) and multiplies at the end:
\begin{equation}
  V_c \;=\; \sum_{\lev} n_{c,\lev}\,\dvol_{\lev},
  \qquad \dvol_{\lev} = \prod_{k=1}^{\dsp} \Delta x_{\lev,k}.
  \label{eq:volume}
\end{equation}
Because uncovered cells partition the domain, \cref{eq:volume} is exact, and
because it is a function of integers it is bit-identical at any decomposition.

\begin{designrule}\label{rule:intout}
Where a scientific output can be derived from integers, derive it from integers.
A quantity that moves in its last bits when the rank count changes is
comparable only to a tolerance, and here that was avoidable rather than
inherent.
\end{designrule}

\section{The sweep}

\texttt{amrSweep} walks the threshold over a linear or logarithmic range and
reports, at each value: the number of connected regions, the total volume, the
largest and median component volume, and the isosurface area. Run on several
candidate surrogates, the tables are directly comparable.

What to look for is a plateau in $N(v)$, or a gap in the size distribution.
Either is a threshold-independent scale and therefore a defensible cell
definition. A featureless continuum is also an answer, and arguably the more
interesting one, because it says that no threshold is defensible and shifts the
question.

The sweep is a coarse sampling of a merge tree, which makes it a free and sharp
regression gate:

\begin{designrule}\label{rule:sweepgate}
When the merge tree is built it must reproduce \texttt{amrSweep}'s $N(v)$
exactly at every sampled threshold.
\end{designrule}

\section{Verification against closed forms}

\texttt{makeTestPlt} writes a cellular field
\begin{equation}
  f(\mathbf{x}) \;=\; \sum_{k=1}^{\dsp} \cos\!\left(\frac{2\pi n\,x_k}{L_k}\right),
  \label{eq:cellfield}
\end{equation}
which on a periodic domain has exactly $n^{\dsp}$ maxima, and saddles between
neighbours at a value that can be written down: moving between adjacent maxima
one term runs from $+1$ to $-1$ while the others stay at $+1$, so the saddle sits
at $\dsp - 2$.

\begin{table}[htbp]
\centering
\caption{Component count against the analytic saddle of \cref{eq:cellfield},
  $n=3$. The maxima lie \emph{on} the domain faces, so recovering exactly
  $n^{\dsp}$ of them is also a test that the periodic wrap is present in the
  adjacency rather than assumed.}
\label{tab:cellcounts}
\begin{tabular}{llll}
\toprule
case & above the saddle & below & analytic saddle \\
\midrule
3D, $n=3$ & 27 components & 1 & exactly $1$ \\
2D, $n=3$ &  9 components & 1 & exactly $0$ \\
\bottomrule
\end{tabular}
\end{table}

A threshold sitting exactly \emph{on} a saddle value gives a genuinely
ill-defined count---the 3D case reads $20$ at $v = 1.000$. That is the
degeneracy persistence exists to quantify, not a defect.

The sphere case checks volume and area against closed forms to $0.06\%$ and
$0.04\%$. Counts and volumes are unchanged across \texttt{max\_grid\_size}
$8/16/32/64$ and across a three-level hierarchy.

\chapter{Verification methodology}
\label{ch:verification}

This chapter is part of the design, not a report on it. The verification
apparatus is the most transferable content in this work, and two of its
principles were learned at the cost of shipped defects.

\section{Analytic references, not plausibility}

Every gate compares against something known in closed form rather than against
a previous run or an eyeball judgement:

\begin{itemize}[leftmargin=*]
\item \textbf{Complex}: the element volume sum against the expected volume---the
  full domain in periodic directions, the domain less one level-0 cell in
  non-periodic ones, where cell-centred data leaves the complex inset by
  $\Delta x/2$ at each wall. Relative error is \emph{exactly zero} across one to
  three levels, refinement ratios $2$ and $4$, periodic, non-periodic and mixed,
  isotropic and anisotropic cells, cubic and elongated domains.
\item \textbf{Surface}: closed with the correct Euler characteristic, and
  second-order convergence of the measure. In 3D the area error on a sphere is
  $-2.799\times10^{-3}$, $-6.985\times10^{-4}$, $-1.745\times10^{-4}$ at
  $32^3$, $64^3$, $128^3$: ratios $4.007$ and $4.003$. In 2D the arc length
  gives $4.01$, $4.00$, $4.00$.
\item \textbf{Carried components}: a linear field interpolated linearly along an
  edge is exact, so a carried $x$ must return equal to the vertex's own $x$
  coordinate. Worst error: exactly zero.
\item \textbf{Segmentation}: \cref{tab:cellcounts}.
\end{itemize}

Compensated summation is used for the volume sum, not out of fastidiousness:
naive summation costs $\sim 5\times10^{-10}$ relative at $2\times 10^6$
elements, and a tolerance loose enough to absorb that would hide a real gap of
the same size. The gate has to stay sharp as the problem grows.

\section{Face conformity, which volume cannot see}
\label{sec:faces}

Volume is additive over any decomposition, so a complex with T-junctions at
coarse--fine interfaces closes \emph{exactly} while being unusable for marching
simplices. It must therefore be measured separately. Each element face is
reduced to its set of distinct identifiers and incidences are counted; on a
periodic domain every non-degenerate face must have incidence exactly two.

Result across one, two and three levels: zero non-manifold faces, zero unmatched
faces. The degenerate faces counted separately---4800 at two levels, 37632 at
three---are the stitch cells' collapsed faces, which have no area and take no
part in sharing. This is also what established the premise the subdivision of
\cref{sec:cone} depends on.

\section{Negative controls}
\label{sec:controls}

Each of these exists because the obvious check was caught passing on
demonstrably broken output. They are retained in the source and should not be
removed.

\begin{table}[htbp]
\centering
\caption{Negative controls and what they reveal about the cheaper checks.}
\label{tab:controls}
\begin{tabular}{p{0.26\textwidth}p{0.28\textwidth}p{0.34\textwidth}}
\toprule
control & what it breaks & what still passes \\
\midrule
\texttt{diagonal\_rule=0} & face-diagonal consistency &
  area---and \emph{better} than the correct rule \\
\texttt{orient\_elements=0} & element winding &
  area \emph{and} closure, both numerically identical \\
\texttt{check\_vertex\_determinism} & element-dependent vertex data &
  everything, in serial \\
improper nesting & stitch with no valid coarse parent &
  nothing: it aborts \\
\bottomrule
\end{tabular}
\end{table}

The first entry is the sharpest lesson available here. Forcing an inconsistent
diagonal produces 10401 boundary edges, 445 non-manifold edges and
$\eulerchi = -2976$, while reporting an area error of $-0.0012$ against the
correct rule's $-0.0016$. \emph{A surface riddled with cracks scored better on
area than a sound one.}

\begin{designrule}\label{rule:notatest}
Area is not a topology test. Closure is not an orientation test. Each is
structurally blind to the defect the next one catches.
\end{designrule}

\section{MPI invariance, and the checksum discipline}
\label{sec:mpi}

Matching counts is weak evidence: two different meshes can share a count. The
apparatus therefore emits order-independent checksums---sums or exclusive-ors of
per-item hashes, invariant to visit order and to which rank holds what, but
sensitive to any difference in membership, connectivity, or the exact bits of
propagated field values.

\begin{designrule}\label{rule:gidkey}
Anything entering a checksum is named by an identifier or a tuple of
identifiers, never by a position in an array.
\end{designrule}

\Cref{rule:gidkey} is stated because it was violated twice. The face census
summed rank-local counts, and the edge parity hashed rank-local vertex indices;
both reported an open surface on a closed one. In both cases the structure was
correct and only the instrument lied.

A related distinction proved necessary. Quantities whose last bits legitimately
depend on the partition---the volume residual, the surface area---must not be
compared for equality, because the per-rank sums are compensated while the
cross-rank reduction is not. Such lines are marked in the output and excluded
from the comparison; the tolerance test carries the same statement in comparable
form. Conditional \emph{presence} is as disqualifying as a drifting value: two
lines printed only at one rank once failed every multi-rank comparison by their
absence.

\section{Guarded regions are untested regions}
\label{sec:guarded}

Code inside \texttt{\#ifdef AMREX\_USE\_MPI} is never compiled by a serial
build. A most-vexing-parse---\texttt{Vector<T> v(std::size\_t(n))} declares a
function---survived every serial test in a gather for exactly that reason.

The response was structural rather than a correction. \texttt{AMRParallel}
reduces the conditional region to the MPI calls themselves: counts, offsets,
resizes and the swap all sit outside the guard and are compiled by every build,
and the collectives are written once instead of repeated at each call site.

\begin{designrule}\label{rule:guard}
Keep guarded regions minimal. Everything that is ordinary code belongs outside
the guard, where every build parses it.
\end{designrule}

\begin{designrule}\label{rule:vexing}
Declare, then resize. Never \texttt{Vector<T> v(std::size\_t(n))}.
\end{designrule}

\section{The pattern in the defects}

Of the defects found by the MPI harness, all but one were in the instrumentation
rather than in the thing instrumented: rank-local face counts, rank-local vertex
indices, a whitespace-split parameter list, a renamed extractor terminator, an
unsorted element list. Only the element-frame vertex position of
\cref{sec:edgekey} was a defect in the code under test.

That is a favourable sign about the design and an unfavourable one about
trusting unfalsified checks, which is the argument for \cref{tab:controls}
existing at all.

\section{Test apparatus}

Three harnesses, all exiting non-zero with a failure count so they can gate:

\begin{itemize}[leftmargin=*]
\item \texttt{Tools/amrComplexMpiTest.sh} --- five sections: rank invariance at
  $1$--$8$ ranks including odd counts; distribution invariance at fixed rank
  count across \texttt{DistributionMapping} strategies, which is the sharper
  test because it varies only \emph{which box lands where}; decomposition
  independence across \texttt{max\_grid\_size}; contour invariance; sweep
  invariance.
\item \texttt{Tools/amrContourMefTest.sh} --- MEF byte-identity across rank
  counts by \texttt{cmp}, with round-trip through \texttt{surfMEFtoDAT}.
\item \texttt{Tools/amrBuildAndTest.sh} --- builds every target in both
  dimensions and runs everything. It verifies each executable exists and is
  executable rather than trusting \texttt{make}'s exit status, because an
  interrupted link once left a non-executable binary that \texttt{make} then
  reported as up to date.
\end{itemize}

Each failure guide states what a given symptom localises to. For MEF, the
byte-layout arithmetic is included, since locating the first differing byte
against the layout is what distinguishes a node-ordering fault from an
element-ordering one with no further instrumentation.

\chapter{The topological target}
\label{ch:topology}

Which topological object to compute is the central design decision, and the
requirements of \cref{ch:problem} settle it. This chapter records the reasoning,
including two conclusions that were wrong in earlier drafts and are corrected
here.

\section{Three candidates}

\begin{table}[htbp]
\centering
\caption{Topological summaries of a scalar field, and what periodicity costs
  each of them.}
\label{tab:candidates}
\begin{tabular}{p{0.20\textwidth}p{0.34\textwidth}p{0.34\textwidth}}
\toprule
object & shape & cost on a periodic domain \\
\midrule
join (merge) tree &
  always a tree, on a torus as much as anywhere &
  nothing \\
contour tree &
  equals the Reeb graph only if the domain is simply connected &
  not available: invalid \\
Reeb graph &
  carries loops &
  loop surgery~\cite{tierny2009loopsurgery} \\
\bottomrule
\end{tabular}
\end{table}

\section{What a Reeb loop is, and is not}

A cycle in the Reeb graph is a one-parameter family of contours that leaves a
level, travels, and returns. Cycles arise from non-trivial $1$-cycles in the
\emph{domain}, not from the shape of any one contour. On a $\dsp$-torus
$\betti{1}(\Tor{\dsp}) = \dsp$, so the Reeb graph of a triply periodic field can
carry up to three independent loops, each meaning a structure that closes on
itself through the wrap. On a non-periodic box $\betti{1} = 0$, the Reeb graph is
a tree, and it is the contour tree.

\paragraph{It does not give the genus of a level set.} A saddle that adds a
handle to a contour without merging or splitting it is a \emph{degree-two} node
in the Reeb graph, and degree-two nodes are conventionally dropped: they sit in
the middle of an arc and are invisible in the output. Handle counts, tunnels and
genus evolution are therefore not readable from a Reeb graph. An earlier draft
claimed otherwise; the correction matters, because it removes the main reason
one might have paid for loop surgery.

The right tool for genus is far cheaper. Extract the surface, split it into
connected components, and compute
\begin{equation}
  \eulerchi \;=\; V - E + F, \qquad g \;=\; \frac{2 - \eulerchi}{2}
  \label{eq:euler}
\end{equation}
per closed orientable component: exact, linear in the element count, no sweep.
Its one prerequisite is that the surface be watertight, which is what
\cref{ch:complex} delivers.

\paragraph{And $\eulerchi$ is the wrong invariant in two dimensions,} where every
closed curve has $\eulerchi = 0$ before and after any event of interest. The 2D
statement is the \emph{number} of boundary components together with their
homology class on the torus: a curve that wraps has no inside, so ``contained
within'' is conditional on being contractible. Winding numbers---net minimum-image
displacement around the loop, divided by the domain length---distinguish the two,
and non-contractible curves necessarily occur in pairs with opposite class,
which is itself a consistency check.

\section{The decision}
\label{sec:decision}

Cells are connected components of a super-level set, and their appearance,
growth and merging is a merge-tree story. Two consequences decide the matter.

\paragraph{Persistence is the robustness measure \ref{req:robust} asked for.}
Rather than choosing a threshold, compute all thresholds at once. Each cell then
carries a birth value $b$, a merge value $m$, and a persistence
\begin{equation}
  p \;=\; |b - m|,
  \label{eq:persistence}
\end{equation}
which is how far the knob must move to make that cell vanish. The reportable
quantity is $N(p)$, the number of cells with persistence above $p$. A plateau in
$N(p)$ is a defensible scale; a featureless continuum says no threshold is
defensible, which is itself the finding and shifts the question rather than
answering it badly.

\paragraph{Periodicity costs nothing.} A join tree is always a tree: components
of $\sublev{v}$ merge monotonically as $v$ decreases and can never form a cycle.
Loop surgery is required for the contour tree and the Reeb graph, neither of
which is needed. The doubly and triply periodic cases of \ref{req:periodic},
with arbitrary face placement, come free---periodicity enters only as adjacency
in the union--find, which \cref{ch:complex} already supplies. The most expensive
of the three candidates was the one not required.

\section{What the merge tree does not cover}

A crescent whose tips join, leaving one contour inside another, is a
$\betti{1}$ event: the region remains one connected component and the merge tree
is blind to it. Concretely, the region is a topological disk with one boundary
curve before the event and an annulus with two after, while
$\betti{0} = 1$ throughout and $\betti{1}$ goes $0 \to 1$. Capturing it requires
the per-component boundary invariants above---boundary-curve counts with homology
class in 2D, \cref{eq:euler} per surface component in 3D---which sit on the
segmentation and are small.

\section{One engine for two comparisons}
\label{sec:overlap}

Matching segmentations between consecutive snapshots by spatial overlap is the
same operation as matching a fuel-consumption segmentation against a temperature
segmentation on the \emph{same} snapshot. Built once it serves both: whether a
cell persists in time (\ref{req:track}), and whether two surrogates identify the
same cells (\ref{req:compare}). The second measures the open question of
\cref{ch:problem} directly---where surrogates agree the structure is
definition-independent, and where they diverge that is differential diffusion
expressed as topology rather than as a contour plot.

Segmentation written as a plotfile is what makes this practical: an integer label
per cell of the hierarchy feeds straight into \texttt{conditionalMean}, so
statistics of any field inside cells defined by any other field is a one-liner.
That is the common context the surrogates need in order to be compared at all,
and it is the presentation used in~\cite{bremer2010burning}.

\paragraph{A hypothesis, not an assertion.} A merge-tree branch that never
becomes an independent maximum---a saddle-separated shoulder of a larger
structure---is structurally distinct from one that is its own high-persistence
branch. If back-supported structures are the former and self-sustaining ones the
latter, then \cref{eq:persistence} measures the degree of independence. This
requires physics judgement and is recorded as a question.

\section{Flame surface thickness}

Broadening under turbulent strain follows from \cref{sec:carried} with no new
machinery: carry $|\nabla f|$ onto the surface as a sampled component, and the
local thickness is
\begin{equation}
  \delta \;\approx\; \frac{\Delta v}{|\nabla f|},
  \label{eq:thickness}
\end{equation}
aggregated over the surface as a distribution rather than a single number.

\chapter{GPU considerations}
\label{ch:gpu}

The implementation described so far is host code. This chapter states that
plainly, separates what would port readily from what would not, and identifies
the one decision that should be taken before the merge tree is written rather
than after.

\section{Current state}

There are no \texttt{ParallelFor} launches in \texttt{AMRComplex},
\texttt{AMRContour} or \texttt{AMRSegment}. Loops are \texttt{AMREX\_LOOP\_3D},
which expands to serial nested loops on the host, and the working data
structures in the hot paths are \texttt{std::map}, \texttt{std::unordered\_map}
and \texttt{std::set}. None of these is usable in device code as written.

This was a deliberate deferral, not an oversight: correctness on coarse--fine
interfaces and periodic domains was the hard part, and establishing it in host
code kept the verification tractable. It is nonetheless a real limitation for a
code base whose production path is GPU.

\section{What ports readily}

\begin{itemize}[leftmargin=*]
\item \textbf{Node classification} (\cref{sec:identity,sec:stitch}). Index
  arithmetic over boxes with no data-dependent control flow. The
  \texttt{AMREX\_LOOP\_3D} bodies map to \texttt{ParallelFor} essentially
  unchanged; \texttt{Node} is already trivially copyable by construction, so it
  can live in device memory.
\item \textbf{Global numbering} (\cref{eq:gid}). A pure function of
  $(\lev,\iv)$ with no lookup and no communication---evaluable per cell on the
  device.
\item \textbf{Element generation} (\cref{sec:ownership}). Ownership is box
  algebra resolved once per grid on the host; the per-cell corner gather and
  validity test that follow are a device loop. A count-then-fill pass with a
  prefix sum replaces \texttt{push\_back}.
\item \textbf{Simplex subdivision and marching} (\cref{sec:cone}). Per-element,
  branch-light, fixed-size output. Well suited to the device.
\item \textbf{Volume and area reductions.} Already reductions; the compensated
  summation of \cref{ch:verification} needs replacing by a device-friendly
  equivalent, and the integer-derived volumes of \cref{eq:volume} are
  unaffected.
\end{itemize}

\section{What requires redesign}

\begin{itemize}[leftmargin=*]
\item \textbf{The vertex cache} (\cref{sec:edgekey}). A \texttt{std::map} from
  edge key to vertex index, with insertion during traversal. The device
  equivalent is to emit one candidate vertex per crossed edge, sort by the
  $64$-bit edge key, and compact---replacing the associative container by a sort
  and a unique. The edge key is already a pair of integers, so this is a change
  of mechanism rather than of design, and \cref{rule:edgefn} is what makes it
  legitimate: because a vertex is a function of its edge alone, duplicates are
  bitwise identical and compaction is exact.
\item \textbf{Union--find} (\cref{ch:segment}). Pointer-chasing with path
  compression over a hash map. A device implementation needs a different
  algorithm---label propagation over a fixed adjacency, or a device-side
  connected-components pass---together with the adjacency materialised as arrays
  rather than discovered by traversal.
\item \textbf{The face and edge censuses} (\cref{sec:faces,sec:orient}). Both
  build \texttt{std::map} tallies. The exclusive-or parities are already
  reductions and port directly; the exact incidence counts would become
  sort-and-count.
\item \textbf{MEF output} (\cref{ch:contour}). Host-side by nature. The
  \texttt{.pvtu} path on the roadmap is the better target and is equally
  host-side, which is appropriate for I/O.
\end{itemize}

\section{A decision to take before the merge tree}

\decide{Whether the merge tree is written as host code, accepting a
device-to-host transfer of the node field per snapshot, or designed for the
device from the start.}

The consideration is that the merge tree is the first component whose cost scales
with the \emph{volume} of data rather than with the area of an extracted surface.
Contouring touches only cells straddling one isovalue; a sweep does so once per
threshold; the merge tree visits every node in sorted order. A global sort of the
node field is the dominant term, and that is a well-suited device operation
whereas the union--find that consumes it is not.

The pragmatic shape, if the device route is chosen, is a device sort feeding a
host tree construction on the resulting order, since the tree itself is small
compared with the field. That should be settled deliberately rather than
inherited from the host-only precedent set here.

A second point worth recording: the verification apparatus of
\cref{ch:verification} is dimension- and backend-agnostic by construction. The
checksums, the analytic references and the negative controls all compare
\emph{outputs}, so a device port can be gated on reproducing the host results
exactly---which is the strongest possible acceptance test for such a port, and is
available at no additional cost.

\chapter{State, roadmap and open decisions}
\label{ch:roadmap}

\section{What exists}

On branch \texttt{amr-complex}, all suites passing in two and three dimensions
at one to eight ranks:

\begin{table}[htbp]
\centering
\caption{Components added.}
\label{tab:files}
\begin{tabular}{p{0.34\textwidth}p{0.58\textwidth}}
\toprule
file & role \\
\midrule
\texttt{Src/AMRComplex.\{H,cpp\}} &
  conforming multilevel complex; identity $(\lev,\iv)$ \\
\texttt{Src/AMRContour.\{H,cpp\}} &
  marching simplices, oriented, carried components, MEF \\
\texttt{Src/AMRSegment.\{H,cpp\}} &
  connected components of a threshold set \\
\texttt{Src/AMRParallel.\{H,cpp\}} &
  the only place MPI calls appear \\
\texttt{Src/amrComplexVerify.cpp} &
  complex gate \\
\texttt{Src/amrContourVerify.cpp} &
  surface gate, negative controls \\
\texttt{Src/amrSweep.cpp} &
  threshold sweep \\
\texttt{Src/makeTestPlt.cpp} &
  synthetic hierarchies with analytic references \\
\texttt{Src/GNUmakefile.AMRComplex} &
  standalone, AMReX only, \texttt{DIM ?= 3} \\
\texttt{Tools/amr*.sh} &
  the three harnesses \\
\bottomrule
\end{tabular}
\end{table}

\section{Roadmap}

\begin{enumerate}[leftmargin=*]
\item \textbf{Representative snapshots.} One plotfile carrying two or three
  candidate surrogates---heat release rate, temperature, an
  intermediate---swept individually, with the $N(v)$ curves laid side by side.
  This answers whether a plateau exists, supplies real timings for
  \texttt{ownedBaseRegions} and the segmentation maps at production size, and
  shows whether the surrogates agree, all before anything is built on an
  assumption about any of the three.
\item \textbf{Segmentation as a plotfile}, and contour component counts.
\item \textbf{Merge tree}, gated on \cref{rule:sweepgate}, and subject to the
  decision of \cref{ch:gpu}.
\item \textbf{Overlap matching} (\cref{sec:overlap}), serving both time tracking
  and cross-surrogate comparison.
\item \textbf{Boundary invariants} for the $\betti{1}$ events.
\end{enumerate}

\section{Open decisions}

\decide{Whether a cell is a region thresholded from above or from below. It may
differ per surrogate, so it remains a runtime choice rather than a design
commitment.}

\decide{Host or device for the merge tree (\cref{ch:gpu}).}

\decide{Wall extrapolation. Cell-centred data leaves the complex inset by
$\Delta x/2$ at each non-periodic wall. Zero normal gradient is a defensible
placeholder and is exactly right for a symmetry, slip or adiabatic wall, but it
is exactly wrong at a Dirichlet wall, and it introduces a prismatic collar whose
area contribution is $P\,\Delta x/2$ for inset-plane perimeter $P$, together
with a spurious curvature ridge standing off each wall at a fixed
$\Delta x/2$. The error is bounded by how obliquely the surface meets the wall,
so the incidence-angle distribution on the inset plane decides whether it
matters at all. Limited linear extrapolation is the recommended default; a
supplied boundary-condition table is the only correct option at a Dirichlet
wall.}

\section{Known limits}

\begin{itemize}[leftmargin=*]
\item MEF output gathers to the I/O processor---the same wall
  \texttt{isosurface.cpp} has. \texttt{.pvtu} is the intended replacement.
\item \texttt{ownedBaseRegions} and the segmentation maps have not been exercised
  at production size.
\item All verification to date uses synthetic fields on generated hierarchies.
  Hierarchies from actual regridding will exercise the proper-nesting assertion
  of \cref{sec:stitch} in a way synthetic ones cannot; if it aborts, that is the
  check working.
\item Host-only, per \cref{ch:gpu}.
\end{itemize}


\appendix
\chapter{Conventions}
\label{app:conventions}

Collected for reference. Each was learned from a defect; breaking one
reintroduces it.

\begin{enumerate}[leftmargin=*]
\item \cref{rule:identity} --- identity is $(\lev,\iv)$, never a coordinate.
  Periodic identification then falls out of \texttt{FillBoundary} with no
  arithmetic, and coarse--fine stitching needs no case table.
\item \cref{rule:gidkey} --- anything entering a checksum is named by an
  identifier or a tuple of identifiers, never by a position in an array.
\item Any ordering written to a file is canonicalised explicitly, never
  inherited from traversal order. Nodes \emph{and} elements
  (\cref{ch:contour}).
\item \cref{rule:guard} --- guarded regions are untested regions. Keep the guard
  around the collectives only.
\item \cref{rule:vexing} --- declare, then resize.
\item \cref{rule:intout} --- derive scientific outputs from integers where
  possible.
\item \cref{rule:edgefn} --- nothing about the element that reached an edge may
  enter what its vertex is.
\item \cref{rule:noparity} --- do not key a subdivision choice on index parity.
\item \cref{rule:notatest} --- area is not a topology test; closure is not an
  orientation test.
\item \cref{rule:sweepgate} --- the merge tree must reproduce the sweep's
  $N(v)$ exactly.
\end{enumerate}

\chapter{Building and testing}
\label{app:build}

\section{This document}

\begin{lstlisting}[style=sh]
cd Docs/design
make
\end{lstlisting}

\section{The code}

\texttt{AMREX\_HOME} defaults to \texttt{../../amrex}, correct for a tree laid
out as \texttt{src/PeleAnalysis} beside \texttt{src/amrex}. The makefile is
independent of \texttt{Make.Analysis} and therefore of PelePhysics, which
\texttt{AMRComplex} does not need.

Everything, both dimensions, with all suites:

\begin{lstlisting}[style=sh]
cd PeleAnalysis
MPIRUN="mpirun --oversubscribe" Tools/amrBuildAndTest.sh
\end{lstlisting}

\texttt{--oversubscribe} is required by OpenMPI once a requested rank count
exceeds the physical core count; MPICH and Intel MPI do not need it. Under Slurm
use \texttt{MPIRUN=srun} from inside an allocation. Honoured:
\texttt{COMP}, \texttt{DIMS}, \texttt{RANKS}, \texttt{JOBS},
\texttt{SKIP\_BUILD}, \texttt{LOGDIR}.

A single target:

\begin{lstlisting}[style=sh]
cd PeleAnalysis/Src
make -f GNUmakefile.AMRComplex USE_MPI=TRUE COMP=llvm DIM=3 EBASE=amrSweep -j8
\end{lstlisting}

\section{Representative invocations}

\begin{lstlisting}[style=sh]
amrComplexVerify3d.llvm.MPI.ex infile=plt00100 is_per="1 1 0" varnames=temp

amrContourVerify3d.llvm.MPI.ex infile=plt00100 is_per="1 1 0" \
    isoCompName=temp isoVal=1500 varnames="Y(OH) HeatRelease" \
    outfile=surf.mef

amrSweep3d.llvm.MPI.ex infile=plt00100 isoCompName=HeatRelease \
    is_per="1 1 0" vlo=1e8 vhi=1e10 nval=40 log=1 outfile=sweep.dat
\end{lstlisting}

Negative controls, which should fail:

\begin{lstlisting}[style=sh]
amrContourVerify3d.llvm.MPI.ex infile=plt00100 is_per="1 1 0" \
    isoCompName=temp isoVal=1500 diagonal_rule=0

amrContourVerify3d.llvm.MPI.ex infile=plt00100 is_per="1 1 0" \
    isoCompName=temp isoVal=1500 orient_elements=0
\end{lstlisting}


\bibliographystyle{plain}
\bibliography{refs}

\end{document}
